% figures/w05-allocation.svg — 배분행렬 B = [1 1; y_p -y_p] 의 기하학적 뜻.
%
%   bash _tools/tikz2svg.sh figures/src/w05-allocation.tex figures/w05-allocation.svg
%
% 사용자 질문 (2026-09-29): "[1 1; 0.395 -0.395] 이게 어떤 의미인지도... 기하학적으로
% 그림 가지고 와서 설명도 해줘봐".
%
% 왼쪽 패널이 답이다. 배 위에서 내려다본 그림 하나에 B 의 네 성분이 전부 들어 있다:
%   1 행 [1  1]        두 추력 모두 **같은 방향으로** 배를 민다 -> 계수가 둘 다 +1
%   2 행 [y_p  -y_p]   좌현은 중심선에서 y_p 만큼 **왼쪽**, 우현은 y_p 만큼 오른쪽에
%                      있으므로 모멘트 팔의 부호가 서로 반대다
% 오른쪽 패널은 그 행렬이 하는 일을 수로 보인다 — 실험 5-7 의 '직진' 행이다.
%
% MSS 데모 demoOtterUSVPathFollowingHeadingControl 의 Allocation 서브시스템에
% Constant [1 1; 0.395 -0.395] 가 그대로 들어 있다. 이 그림은 그 상수의 그림이다.
%
% 좌표는 전부 여기에 손으로 적혀 있고 산술이 없다 — 선체는 대칭이므로 y_p 만 정하면
% 나머지가 따라온다 (y_p = 0.395 m 를 화면에서 1.6 cm 로 잡았다, 축척 4.05 cm/m).
\documentclass[border=5pt]{standalone}
% gnc-style.tex — 이 볼트의 TikZ 그림이 공유하는 스타일.
%
% 저널 그림 규약을 스타일 하나로 굳혀 둔다. 그림마다 선 굵기나 화살촉을
% 다시 정하지 않는다 — 그것이 그림이 제각각으로 보이는 원인이다.
%
% 쓰는 법:  % gnc-style.tex — 이 볼트의 TikZ 그림이 공유하는 스타일.
%
% 저널 그림 규약을 스타일 하나로 굳혀 둔다. 그림마다 선 굵기나 화살촉을
% 다시 정하지 않는다 — 그것이 그림이 제각각으로 보이는 원인이다.
%
% 쓰는 법:  % gnc-style.tex — 이 볼트의 TikZ 그림이 공유하는 스타일.
%
% 저널 그림 규약을 스타일 하나로 굳혀 둔다. 그림마다 선 굵기나 화살촉을
% 다시 정하지 않는다 — 그것이 그림이 제각각으로 보이는 원인이다.
%
% 쓰는 법:  \input{gnc-style}  (figures/src/ 안에서)

\usepackage{tikz}
\usepackage{amsmath}
\usepackage{xcolor}
\usetikzlibrary{arrows.meta, positioning, calc, fit, backgrounds, decorations.pathmorphing}

% ---- 색 --------------------------------------------------------------------
% 색은 강조 하나에만 쓴다. 나머지는 검정.
\definecolor{ink}{HTML}{1A1A1A}
\definecolor{accent}{HTML}{7C3AED}
\definecolor{warn}{HTML}{B91C1C}
\definecolor{soft}{HTML}{666666}

% 기하 그림(좌표계·유도)에서만 쓰는 넷. 블록선도는 위의 흑백 + accent 로 충분하다.
\definecolor{pathc}{HTML}{0D9488}   % 계획 경로 · 항적
\definecolor{errc}{HTML}{C2410C}    % 오차 (cross-track 등)
\definecolor{alng}{HTML}{B45309}    % along-track
\definecolor{flow}{HTML}{0369A1}    % 조류 · 유속

% ---- 선과 화살촉 -----------------------------------------------------------
\tikzset{
  % 신호선. 화살촉은 작고 뾰족한 것 하나뿐이다.
  sig/.style   = {ink, line width=0.5pt, -{Stealth[length=4pt, width=3pt]}},
  wire/.style  = {ink, line width=0.5pt},
  % 강조 경로 — 파선으로 구분한다. 흑백 인쇄에서도 살아야 하기 때문이다.
  asig/.style  = {accent, line width=0.5pt, densely dashed,
                  -{Stealth[length=4pt, width=3pt]}},
  awire/.style = {accent, line width=0.5pt, densely dashed},
  % 블록. 흰 바탕, 직각, 검은 테두리.
  blk/.style   = {draw=ink, line width=0.55pt, fill=white,
                  minimum width=17mm, minimum height=8mm, inner sep=2pt, font=\small},
  % 합산점
  sum/.style   = {draw=ink, line width=0.55pt, fill=white, circle,
                  inner sep=0pt, minimum size=4.6mm, font=\scriptsize},
  asum/.style  = {draw=accent, line width=0.55pt, fill=white, circle,
                  inner sep=0pt, minimum size=4.6mm, font=\scriptsize},
  % 분기점
  dot/.style   = {ink, fill=ink, circle, inner sep=0pt, minimum size=1.5mm},
  % 신호 이름 — 선 위에 작게
  sname/.style = {font=\small, inner sep=1.5pt},
  % 그림 안의 짧은 설명. 그림 안의 글은 최소로 둔다
  note/.style  = {font=\scriptsize, soft, inner sep=1.5pt},
  anote/.style = {font=\scriptsize, accent, inner sep=1.5pt},
  % 축
  axis/.style  = {ink, line width=0.5pt},
  tick/.style  = {font=\scriptsize, soft},
}

% 캡션 한 줄 — 그림 아래 가는 선과 함께
\newcommand{\gnccaption}[2]{%
  \node[anchor=north west, text width=#1, align=left, font=\scriptsize, ink]
        at (current bounding box.south west) {\vspace{2pt}#2};}
  (figures/src/ 안에서)

\usepackage{tikz}
\usepackage{amsmath}
\usepackage{xcolor}
\usetikzlibrary{arrows.meta, positioning, calc, fit, backgrounds, decorations.pathmorphing}

% ---- 색 --------------------------------------------------------------------
% 색은 강조 하나에만 쓴다. 나머지는 검정.
\definecolor{ink}{HTML}{1A1A1A}
\definecolor{accent}{HTML}{7C3AED}
\definecolor{warn}{HTML}{B91C1C}
\definecolor{soft}{HTML}{666666}

% 기하 그림(좌표계·유도)에서만 쓰는 넷. 블록선도는 위의 흑백 + accent 로 충분하다.
\definecolor{pathc}{HTML}{0D9488}   % 계획 경로 · 항적
\definecolor{errc}{HTML}{C2410C}    % 오차 (cross-track 등)
\definecolor{alng}{HTML}{B45309}    % along-track
\definecolor{flow}{HTML}{0369A1}    % 조류 · 유속

% ---- 선과 화살촉 -----------------------------------------------------------
\tikzset{
  % 신호선. 화살촉은 작고 뾰족한 것 하나뿐이다.
  sig/.style   = {ink, line width=0.5pt, -{Stealth[length=4pt, width=3pt]}},
  wire/.style  = {ink, line width=0.5pt},
  % 강조 경로 — 파선으로 구분한다. 흑백 인쇄에서도 살아야 하기 때문이다.
  asig/.style  = {accent, line width=0.5pt, densely dashed,
                  -{Stealth[length=4pt, width=3pt]}},
  awire/.style = {accent, line width=0.5pt, densely dashed},
  % 블록. 흰 바탕, 직각, 검은 테두리.
  blk/.style   = {draw=ink, line width=0.55pt, fill=white,
                  minimum width=17mm, minimum height=8mm, inner sep=2pt, font=\small},
  % 합산점
  sum/.style   = {draw=ink, line width=0.55pt, fill=white, circle,
                  inner sep=0pt, minimum size=4.6mm, font=\scriptsize},
  asum/.style  = {draw=accent, line width=0.55pt, fill=white, circle,
                  inner sep=0pt, minimum size=4.6mm, font=\scriptsize},
  % 분기점
  dot/.style   = {ink, fill=ink, circle, inner sep=0pt, minimum size=1.5mm},
  % 신호 이름 — 선 위에 작게
  sname/.style = {font=\small, inner sep=1.5pt},
  % 그림 안의 짧은 설명. 그림 안의 글은 최소로 둔다
  note/.style  = {font=\scriptsize, soft, inner sep=1.5pt},
  anote/.style = {font=\scriptsize, accent, inner sep=1.5pt},
  % 축
  axis/.style  = {ink, line width=0.5pt},
  tick/.style  = {font=\scriptsize, soft},
}

% 캡션 한 줄 — 그림 아래 가는 선과 함께
\newcommand{\gnccaption}[2]{%
  \node[anchor=north west, text width=#1, align=left, font=\scriptsize, ink]
        at (current bounding box.south west) {\vspace{2pt}#2};}
  (figures/src/ 안에서)

\usepackage{tikz}
\usepackage{amsmath}
\usepackage{xcolor}
\usetikzlibrary{arrows.meta, positioning, calc, fit, backgrounds, decorations.pathmorphing}

% ---- 색 --------------------------------------------------------------------
% 색은 강조 하나에만 쓴다. 나머지는 검정.
\definecolor{ink}{HTML}{1A1A1A}
\definecolor{accent}{HTML}{7C3AED}
\definecolor{warn}{HTML}{B91C1C}
\definecolor{soft}{HTML}{666666}

% 기하 그림(좌표계·유도)에서만 쓰는 넷. 블록선도는 위의 흑백 + accent 로 충분하다.
\definecolor{pathc}{HTML}{0D9488}   % 계획 경로 · 항적
\definecolor{errc}{HTML}{C2410C}    % 오차 (cross-track 등)
\definecolor{alng}{HTML}{B45309}    % along-track
\definecolor{flow}{HTML}{0369A1}    % 조류 · 유속

% ---- 선과 화살촉 -----------------------------------------------------------
\tikzset{
  % 신호선. 화살촉은 작고 뾰족한 것 하나뿐이다.
  sig/.style   = {ink, line width=0.5pt, -{Stealth[length=4pt, width=3pt]}},
  wire/.style  = {ink, line width=0.5pt},
  % 강조 경로 — 파선으로 구분한다. 흑백 인쇄에서도 살아야 하기 때문이다.
  asig/.style  = {accent, line width=0.5pt, densely dashed,
                  -{Stealth[length=4pt, width=3pt]}},
  awire/.style = {accent, line width=0.5pt, densely dashed},
  % 블록. 흰 바탕, 직각, 검은 테두리.
  blk/.style   = {draw=ink, line width=0.55pt, fill=white,
                  minimum width=17mm, minimum height=8mm, inner sep=2pt, font=\small},
  % 합산점
  sum/.style   = {draw=ink, line width=0.55pt, fill=white, circle,
                  inner sep=0pt, minimum size=4.6mm, font=\scriptsize},
  asum/.style  = {draw=accent, line width=0.55pt, fill=white, circle,
                  inner sep=0pt, minimum size=4.6mm, font=\scriptsize},
  % 분기점
  dot/.style   = {ink, fill=ink, circle, inner sep=0pt, minimum size=1.5mm},
  % 신호 이름 — 선 위에 작게
  sname/.style = {font=\small, inner sep=1.5pt},
  % 그림 안의 짧은 설명. 그림 안의 글은 최소로 둔다
  note/.style  = {font=\scriptsize, soft, inner sep=1.5pt},
  anote/.style = {font=\scriptsize, accent, inner sep=1.5pt},
  % 축
  axis/.style  = {ink, line width=0.5pt},
  tick/.style  = {font=\scriptsize, soft},
}

% 캡션 한 줄 — 그림 아래 가는 선과 함께
\newcommand{\gnccaption}[2]{%
  \node[anchor=north west, text width=#1, align=left, font=\scriptsize, ink]
        at (current bounding box.south west) {\vspace{2pt}#2};}


\definecolor{knownc}{HTML}{1D4ED8}

\tikzset{
  hull/.style    = {ink, line width=0.9pt, fill=black!4},
  pont/.style    = {ink, line width=0.7pt, fill=pathc!12},
  thrust/.style  = {errc, line width=1.4pt, -{Stealth[length=5.0pt,width=3.6pt]}},
  sumline/.style = {alng, line width=1.4pt, -{Stealth[length=5.0pt,width=3.6pt]}},
  armline/.style = {knownc, line width=0.8pt, {Stealth[length=3.6pt,width=2.6pt]}-{Stealth[length=3.6pt,width=2.6pt]}},
  ctr/.style     = {soft, line width=0.4pt, densely dashed},
  spin/.style    = {alng, line width=1.1pt, -{Stealth[length=4.4pt,width=3.2pt]}},
  lbl/.style     = {font=\scriptsize},
  known/.style   = {draw=knownc, line width=0.6pt, fill=white, inner sep=3.2pt,
                    rounded corners=1pt, font=\small},
  want/.style    = {draw=accent, line width=0.9pt, fill=white, inner sep=3.4pt,
                    rounded corners=1pt, font=\small},
  capt/.style    = {font=\scriptsize\bfseries, warn},
}

\begin{document}
\begin{tikzpicture}[x=1cm, y=1cm]

% =====================================================================
% 왼쪽: 배를 위에서 본 그림. B 의 네 성분이 전부 여기 있다.
% =====================================================================
\begin{scope}[shift={(0,0)}]

  % 선체 중심선 / the centreline
  \draw[ctr] (0,-2.7) -- (0,2.9);
  \node[lbl, soft, anchor=south] at (0,2.9) {centreline};

  % 두 폰툰. y_p = 0.395 m 를 1.6 cm 로 그린다 (좌현 x=-1.6, 우현 x=+1.6)
  \draw[pont] (-2.05,-2.1) rectangle (-1.15,1.5);
  \draw[pont] ( 1.15,-2.1) rectangle ( 2.05,1.5);
  % 갑판 / the deck
  \draw[hull] (-2.05,0.4) rectangle (2.05,1.5);

  % 선체 좌표축: x 는 앞, y 는 우현 / body axes: x forward, y to starboard
  \draw[ink, line width=0.7pt, -{Stealth[length=4.2pt,width=3.0pt]}] (0,0) -- (0,2.55);
  \node[lbl, anchor=west] at (0.08,2.45) {$x_b$ (forward)};
  \draw[ink, line width=0.7pt, -{Stealth[length=4.2pt,width=3.0pt]}] (0,0) -- (2.75,0);
  \node[lbl, anchor=south] at (2.6,0.08) {$y_b$ (starboard)};
  \fill[ink] (0,0) circle (1.3pt);
  \node[lbl, anchor=north east] at (-0.06,-0.06) {origin};

  % 모멘트 팔 / the moment arms
  \draw[armline] (-1.6,-1.35) -- (0,-1.35);
  \draw[armline] ( 0,-1.35)   -- (1.6,-1.35);
  \node[lbl, knownc, anchor=north] at (-0.8,-1.42) {$y_p$};
  \node[lbl, knownc, anchor=north] at ( 0.8,-1.42) {$y_p$};
  \node[lbl, knownc, anchor=north, align=center] at (0,-1.95)
       {$y_p = 0.395$ m\\[-1pt]half the pontoon spacing};

  % 두 추력. 둘 다 앞으로 / both thrusts point forward
  \draw[thrust] (-1.6,1.5) -- (-1.6,3.05);
  \draw[thrust] ( 1.6,1.5) -- ( 1.6,2.55);
  \node[lbl, errc, anchor=south] at (-1.6,3.08) {$T_L$};
  \node[lbl, errc, anchor=south] at ( 1.6,2.58) {$T_R$};
  \node[lbl, errc, anchor=east, align=right] at (-1.72,2.3)
       {port pushes\\[-2pt]harder};

  % 그 결과의 회전 / the turn it produces
  %  회전 화살표는 y_b 축 **아래** 사분면에 둔다. 위에 두면 'starboard' 이름표와
  %  겹쳐 찍힌다 — 렌더해서 보고 옮겼다 (CLAUDE.md §4 규칙 6).
  \draw[spin] (0.35,-0.94) arc (-70:-18:1.0);
  \node[lbl, alng, anchor=west] at (1.10,-0.55) {$N>0$};

  \node[capt, anchor=north, align=center] at (0,-2.55)
       {one picture, four numbers};
\end{scope}

% =====================================================================
% 가운데: 그 그림에서 읽은 두 줄, 그리고 행렬
% =====================================================================
\begin{scope}[shift={(4.9,0)}]
  \node[known, anchor=north west, align=left] at (0,2.95)
    {\textbf{row 1} — both push forward\\[2pt]
     $X = (+1)\,T_L + (+1)\,T_R$};
  \node[known, anchor=north west, align=left] at (0,1.55)
    {\textbf{row 2} — opposite arms\\[2pt]
     $N = (+y_p)\,T_L + (-y_p)\,T_R$};
  \node[want, anchor=north west, align=left] at (0,-0.05)
    {$\begin{bmatrix} X \\ N \end{bmatrix}
      = \underbrace{\begin{bmatrix} 1 & 1 \\ y_p & -y_p \end{bmatrix}}_{\mathbf{B}}
        \begin{bmatrix} T_L \\ T_R \end{bmatrix}$};
  \node[anchor=north west, align=left, font=\scriptsize] at (0,-1.55)
    {$\det \mathbf{B} = -2 y_p \neq 0$, so\\[2pt]
     $\begin{bmatrix} T_L \\ T_R \end{bmatrix}
       = \mathbf{B}^{-1} \begin{bmatrix} X \\ N \end{bmatrix}
       = \begin{bmatrix} X/2 + N/(2y_p) \\ X/2 - N/(2y_p) \end{bmatrix}$};
  \node[capt, anchor=north, align=center] at (2.4,-2.55)
       {read off the picture, then invert};
\end{scope}

% =====================================================================
% 오른쪽: 실험 5-7 의 '직진' 행을 그대로 / the 'straight leg' row of Experiment 5-7
% =====================================================================
\begin{scope}[shift={(11.0,0)}]
  \node[anchor=north west, align=left, font=\scriptsize] at (0,2.95)
    {\textbf{Experiment 5-7, straight leg}\\[3pt]
     asked for\\[1pt]
     \quad $X = 77.59$ N,\; $N = -0.50$ N\,m\\[4pt]
     solve $\mathbf{B}\,\mathbf{T} = [X;\,N]$\\[1pt]
     \quad $T_L = 38.16$ N,\; $T_R = 39.43$ N\\[4pt]
     then $T = k\,n\lvert n\rvert$\\[1pt]
     \quad $n_L = 58.68$,\; $n_R = 59.65$ rad/s};
  \node[anchor=north west, align=left, font=\scriptsize, soft] at (0,-0.35)
    {\textbf{hardest turn}\\[3pt]
     \quad $X = 120.00$ N,\; $N = -47.15$ N\,m\\[3pt]
     \quad $T_L = 0.32$ N,\; $T_R = 119.68$ N\\[3pt]
     \quad $n_R = 103.93 = n_{\max}$ exactly};
  \node[capt, anchor=north, align=center] at (2.3,-2.55)
       {the numbers the model logs};
\end{scope}

\end{tikzpicture}
\end{document}
